\chapter{结论与展望}
	\subsection{主要结论}
	（1）通过HITRAN数据库line-by-line仿真确定2.7~$\mu$m（3696~cm$^{-1}$）为最佳实验波段，温度灵敏度2.5\%/K，谱线对$\Delta E''\approx660$~cm$^{-1}$保障了低温区间（200--300~K）的测温精度。三点法Abel反演在3\%噪声下误差4.2\%，确定为实测数据核心算法。
	（2）搭建了完整的TDLAS层析诊断系统：聚焦光路将光斑从2.5~mm缩至0.3~mm（信号增强3倍），独立接地使探测器噪声从50~mV降至10~mV（下降80\%），气密吹扫使环境CO$_2$干扰降低96\%。Red Pitaya自触发同步方案解决了相位抖动问题，500次软件平均后SNR提升至30。
	（3）获取3平面$\times$6点共18组层析光谱数据，三点法Abel反演给出射流半宽约1.6~mm。双线强度比法反演温度：射流中心268~K，径向递增至305~K，与Fluent仿真（中心262~K）偏差2.3\%，与Wu等（2021）瑞利散射结果（255-265~K）吻合，验证了TDLAS+Abel反演方法的有效性。
	（4）基于准一维相变模型，以实测温压场为输入，预测凝结起始位置约$z=0.8$~mm，成核峰值位于$z\approx1.5$~mm，马赫盘后团簇蒸发趋势与文献一致。
	\subsection{创新点}
	（1）将TDLAS层析技术应用于欠膨胀CO$_2$超声速射流温度定量测量，首次获取射流径向温度分布（中心268~K→径向305~K）。
	（2）建立“CFD仿真+Abel反演+TDLAS实验”三位一体验证框架，形成从仿真到实验的完整闭环验证体系。
	（3）搭建集成聚焦光路（0.3~mm）、独立接地（噪声↓80\%）、气密吹扫（干扰↓96\%）的高信噪比诊断系统，SNR提升至30。
	\subsection{下一步工作}
	\textbf{短期（1-2个月）}：
	（1）修正FSR标定代码（峰值检测算法改进），提升波长标定精度；
	（2）确定FPGA硬件平均方案（重建Red Pitaya PS端或采用备用方案），缩短单点采集时间。
	\textbf{中期（3-4个月）}：
	（3）扩大扫描至5平面$\times$10点，获取更完整温度场；
	（4）引入瑞利散射或团簇成像，直接探测相变信号；
	（5）优化相变模型数值方法（调整网格密度、松弛因子，逐步加入非定常项），解决小喉部面积工况下的收敛问题。
	\textbf{长期（5-6个月）}：
	（6）基于实验数据校准CNT模型参数（Tolman修正$D$值、非平衡修正因子）；
	（7）校准后模型嵌入CFD框架验证；
	（8）撰写小论文投稿。
	% ==================== 参考文献 ====================

